% Generated by code/interviews-1/03-speed-drills/python/gen_drill.py. Do not edit. French copy by draft/pilot-quant-fr/ivgen_fr.py.
\ivfamily{Additionner et soustraire de grands nombres}{1, 24, 37, 38, 44, 45, 59, 65, 68, 72, 74}{Arrondir un nombre puis corriger.}
\ivfamily{Deux chiffres fois deux chiffres}{8, 16, 18, 21, 26, 29, 30, 35, 46, 50, 58, 66, 73}{Décomposer un facteur, ou utiliser la différence de deux carrés si les facteurs sont proches.}
\ivfamily{Carrés}{7, 14, 39, 47, 52, 63}{$n^2 = (n - d)(n + d) + d^2$.}
\ivfamily{Autour d'un nombre rond}{17, 23, 25, 28, 60}{$(100 + d)(100 - d) = 10000 - d^2$.}
\ivfamily{Pourcentages}{4, 13, 15, 22, 27, 33, 40, 41, 42, 49, 53, 56}{Construire à partir de $10\%$ et de $1\%$. Pour une variation, multiplier par $1 + r$.}
\ivfamily{Fractions et décimaux}{9, 10, 11, 12, 51, 71, 76, 79, 80}{Connaître les huitièmes, seizièmes et quarantièmes. Simplifier avant de multiplier.}
\ivfamily{Conversions}{3, 61, 64, 67}{Ramener une fraction à 100 pour obtenir un pourcentage. Lire un décimal sur une puissance de 10 pour obtenir une fraction.}
\ivfamily{Puissances et racines}{6, 19, 20, 31, 36, 54, 55, 57, 62, 77}{Connaître les puissances de 2 jusqu'à $2^{16}$ et les carrés jusqu'à 29.}
\ivfamily{Coefficients binomiaux}{34, 43, 48, 70, 75, 78}{Simplifier avant de multiplier.}
\ivfamily{Algèbre}{2, 5, 32, 69}{Défaire chaque étape mentalement.}
